\documentclass[11pt]{article}
\usepackage[margin=1in]{geometry}
\usepackage{booktabs}
\usepackage{hyperref}
\usepackage{amsmath}
\usepackage{url}

\title{Memo 03: MP1 Mass-Conservation Check and Updated Diagnosis}
\author{FVCOM-MP Delaware Test Cases}
\date{May 1, 2026}

\begin{document}
\sloppy
\maketitle

\section*{Purpose}
This memo updates the MP1 mass-conservation diagnosis after rerunning cases
\texttt{b} and \texttt{c} with the same sediment input now used by case
\texttt{a}.  The present build still bypasses the sediment horizontal
\texttt{adv\_scal\_backward} and \texttt{fct\_sed} block in \texttt{mod\_sed.F}.
The new result is a clean plastic split-path test: cases \texttt{b} and
\texttt{c} now evolve almost identically, track case \texttt{a} for about three
days, and then diverge together onto an oscillatory, lower-mass branch.

\section*{Archived Analysis Outputs}
Six mass-conservation snapshots are now preserved under
\path{FVCOM_MP_test_run/PYTHON/output/mass_conservation/archive/}.

\begin{table}[h]
\centering
\small
\begin{tabular}{lp{0.34\linewidth}p{0.43\linewidth}}
\toprule
Version & Archive label & Meaning \\
\midrule
\texttt{v001} & \texttt{v001\_before\_case\_a\_sediment\_on} &
old comparison with case \texttt{a} using \texttt{SEDIMENT\_MODEL=F} \\
\texttt{v002} & \texttt{v002\_case\_a\_sediment\_on\_full10d} &
case \texttt{a} with \texttt{SEDIMENT\_MODEL=T} and the original sediment BBL
settings \\
\texttt{v003} & \texttt{v003\_case\_a\_sediment\_on\_no\_bbl} &
case-\texttt{a} diagnostic with \texttt{SEDIMENT\_MODEL=T} and sediment BBL
roughness switches disabled \\
\texttt{v004} & \texttt{v004\_sed\_start\_delayed} &
case \texttt{a} with sediment initialized but active \texttt{Advance\_Sed}
delayed beyond the run \\
\texttt{v005} & \texttt{v005\_sed\_no\_adv\_fct} &
case \texttt{a} with active sediment runtime, but sediment
\texttt{adv\_scal\_backward}/\texttt{fct\_sed} bypassed \\
\texttt{v006} & \texttt{v006\_b\_c\_mirrored\_sediment} &
cases \texttt{b}/\texttt{c} rerun with the same sediment input as case
\texttt{a}; all cases use \texttt{mp\_sediment\_a\_no\_adv.inp} \\
\bottomrule
\end{tabular}
\end{table}

The live mass-conservation output folder has also been regenerated from the
\texttt{v006} state.  The archived files include \texttt{mass\_summary.csv},
\texttt{mass\_timeseries\_long.csv}, and all current comparison plots.

\section*{Case-\texttt{a} Evidence Trail}
All archived comparisons cover the same FVCOM model-day window,
\(58485.0\)--\(58494.957031\), or comparison day \(0.00\)--\(9.96\).

\begin{table}[h]
\centering
\small
\begin{tabular}{lp{0.45\linewidth}rr}
\toprule
Version & Case-\texttt{a} sediment treatment & Final total (kg) &
Final vs. \texttt{v001} (kg) \\
\midrule
\texttt{v001} & sediment model off & 1001.25 & 0.00 \\
\texttt{v002} & sediment on, original BBL settings & 582.41 & -418.85 \\
\texttt{v003} & sediment on, BBL disabled & 573.28 & -427.97 \\
\texttt{v004} & sediment initialized, \texttt{Advance\_Sed} delayed & 1001.25 & 0.00 \\
\texttt{v005} & sediment on, \texttt{adv\_scal\_backward}/\texttt{fct\_sed}
bypassed & 994.20 & -7.05 \\
\bottomrule
\end{tabular}
\end{table}

The two low-mass runs are \texttt{v002} and \texttt{v003}, both of which execute
active sediment horizontal advection and FCT limiting.  The two high-mass
controls are \texttt{v001} and \texttt{v004}: either sediment is absent, or
sediment setup is present but active \texttt{Advance\_Sed} returns before
runtime sediment transport.  The new \texttt{v005} result is the key isolation:
it keeps active sediment runtime and starts sediment immediately
(\texttt{SED\_START = 1}), but bypasses the sediment horizontal
advection/FCT block.  Case \texttt{a} then recovers to \(994.20~\mathrm{kg}\),
within \(7.05~\mathrm{kg}\) of the old \texttt{v001}/\texttt{v004} level.

\section*{Current \texttt{v006} Output Summary}
\begin{table}[h]
\centering
\small
\begin{tabular}{lp{0.38\linewidth}rrr}
\toprule
Case & Current output state & Final water (kg) & Final bed (kg) &
Final total (kg) \\
\midrule
\texttt{a} & rerun with sediment advection/FCT bypassed & 994.20 & 0.00 &
994.20 \\
\texttt{b} & mirrored sediment input, \texttt{PLAST\_FLOC=T} & 347.56 & 0.00 &
347.56 \\
\texttt{c} & mirrored sediment input, \texttt{PLAST\_FLOC=T} & 348.52 & 0.00 &
348.52 \\
\bottomrule
\end{tabular}
\end{table}

The hard concentration cap remains inactive.  The largest saved MP1
concentration in the \texttt{v006} live analysis is
\(4.37\times10^{-4}~\mathrm{kg\,m^{-3}}\) for case \texttt{a},
\(2.38\times10^{-4}~\mathrm{kg\,m^{-3}}\) for case \texttt{b}, and
\(2.38\times10^{-4}~\mathrm{kg\,m^{-3}}\) for case \texttt{c}; all cap counters
are zero.  The saved split fields are exactly internally consistent in the new
outputs: \(\texttt{mp1\_agg}+\texttt{mp1\_dis}=\texttt{mp1}\) to the saved
precision for both \texttt{b} and \texttt{c}.  The aggregated component remains
zero throughout both runs; all water-column mass is in the disaggregated
\texttt{mp1\_dis} component.

\section*{\texttt{v006} Divergence Timing}
Cases \texttt{b} and \texttt{c} are now effectively the same experiment.  Their
maximum total-mass separation over the full run is only \(1.68~\mathrm{kg}\),
and their final totals differ by \(0.96~\mathrm{kg}\).  This confirms that the
old \texttt{b}/\texttt{c} separation was driven by the different sediment files,
not by a separate plastic behavior.

\begin{table}[h]
\centering
\small
\begin{tabular}{rrrrr}
\toprule
Comparison day & Case \texttt{a} (kg) & Case \texttt{b} (kg) &
Case \texttt{c} (kg) & \texttt{a}-\texttt{b} (kg) \\
\midrule
1.00 & 20.13 & 20.13 & 20.13 & 0.00 \\
2.00 & 73.25 & 73.25 & 73.25 & 0.00 \\
3.00 & 147.34 & 141.87 & 141.87 & 5.48 \\
4.00 & 237.41 & 181.30 & 181.32 & 56.11 \\
5.00 & 371.85 & 222.70 & 223.07 & 149.15 \\
9.96 & 994.20 & 347.56 & 348.52 & 646.64 \\
\bottomrule
\end{tabular}
\end{table}

The first nontrivial \texttt{a}/\texttt{b} split occurs just before day 3:
\(|a-b|>1~\mathrm{kg}\) at day \(2.957031\), \(|a-b|>5~\mathrm{kg}\) at day
\(3.000000\), and \(|a-b|>50~\mathrm{kg}\) at day \(4.000000\).  After day 3,
case \texttt{a} has no sign changes in its step-to-step total-mass tendency,
while cases \texttt{b} and \texttt{c} each have 28 sign changes.  Thus the new
failure mode is not a smooth offset; it is an oscillatory split-branch response
that starts once the plume has evolved for about three days.

\section*{Relevant Call Trace}
\begin{enumerate}
\item \textbf{Setup order.} In \texttt{fvcom.F:643--651},
\texttt{SETUP\_SED} is called before \texttt{SETUP\_PLAST}.

\item \textbf{Timestep order.} In \texttt{internal\_step.F:1092--1137},
\texttt{ADVANCE\_SED} is called before \texttt{ADVANCE\_PLAST}.  Plastic
therefore sees whatever state sediment leaves behind each internal step.

\item \textbf{Sediment startup guard.} In \texttt{mod\_sed.F:883},
\texttt{Advance\_Sed} returns immediately when \texttt{iint < sed\_start}.  The
\texttt{v004} test used this to keep sediment setup active while removing
active sediment runtime.

\item \textbf{Bypassed block in the current build.} In
\texttt{mod\_sed.F:1016--1033}, the sediment
\texttt{adv\_scal\_backward} call and \texttt{fct\_sed} limiter are commented,
while \texttt{sed(i)\%cnew = sed(i)\%conc} and
\texttt{sed(i)\%cflx = 0.0\_sp} keep the rest of \texttt{Advance\_Sed} moving.

\item \textbf{Shared scalar machinery.} \texttt{Adv\_Scal\_Backward} begins at
\texttt{mod\_scal.F:648}.  During that routine, bottom scalar-gradient arrays
\texttt{PFPXB}/\texttt{PFPYB} are written at
\texttt{mod\_scal.F:911--912}.  This remains a concrete shared-state candidate,
although the current test only proves that the sediment
\texttt{adv\_scal\_backward}/\texttt{fct\_sed} stage is implicated, not which
line inside the stage is responsible.

\item \textbf{Split plastic branch.} In the \texttt{PLAST\_FLOC=T} path,
\texttt{mod\_plast.F:1166--1186} transports \texttt{conc\_a} and
\texttt{conc\_d} with separate \texttt{adv\_scal} calls.  It then applies
vertical diffusion separately at \texttt{mod\_plast.F:1203--1204}, applies
split OBC/source handling at \texttt{mod\_plast.F:1543--1605}, clamps
\texttt{cnew\_a}/\texttt{cnew\_d} at \texttt{mod\_plast.F:1657--1660}, and
finally recombines at \texttt{mod\_plast.F:1709--1716}.  In the new
\texttt{v006} runs, \texttt{SOURCE\_AG\_RATE=0} and
\texttt{OBC\_AG\_RATE=0}, so this branch should reduce to a disaggregated-only
version of the no-floc case.

\item \textbf{Symmetry check.} At the direct call-signature level,
\texttt{conc\_d} is close to the no-floc \texttt{conc} path:
\texttt{adv\_scal(conc\_d,cnew\_d,\ldots)} uses
\(\texttt{cdis}\times(1-\texttt{Arate\_s0})\), which equals
\texttt{cdis} when \texttt{SOURCE\_AG\_RATE=0}; the matching no-floc call is
\texttt{adv\_scal(conc,cnew,\ldots)} at \texttt{mod\_plast.F:1261--1267}.
However, the full operator path is not directly symmetric because the split
branch still executes floc-rate setup, split deposition/erosion code, separate
zero-field aggregate transport, split OBC/source calls, split clamping, split
MPI exchanges, and a later recombination step.
\end{enumerate}

\section*{Diagnosis}
\begin{enumerate}
\item \textbf{The problem is no longer just ``sediment runtime on''.}
\texttt{v002} and \texttt{v003} showed that active sediment runtime could move
case-\texttt{a} final domain MP1 mass from about \(1001~\mathrm{kg}\) down to
\(573\)--\(582~\mathrm{kg}\).  \texttt{v004} then showed this was not caused by
sediment setup, sediment namelist presence, or the sediment-enabled river
forcing path, because delaying active \texttt{Advance\_Sed} recovered the old
mass curve.

\item \textbf{The new smoking gun is the sediment horizontal
advection/FCT stage.} \texttt{v005} keeps active sediment runtime, starts
sediment immediately, and keeps the non-advection parts of
\texttt{Advance\_Sed} active.  Closing only the sediment
\texttt{adv\_scal\_backward}/\texttt{fct\_sed} block recovers the
\texttt{v001}/\texttt{v004}-like case-\texttt{a} result.  That makes this block
the leading source of the unwanted sediment/plastic interaction.

\item \textbf{\texttt{v006} moves the active problem to the plastic split
path.} With the sediment input now mirrored, cases \texttt{b} and \texttt{c}
are nearly identical.  The old \texttt{b}/\texttt{c} difference was therefore a
sediment-input difference, not a separate plastic mechanism.  The remaining
large \texttt{a} versus \texttt{b}/\texttt{c} split is now controlled by
\texttt{PLAST\_FLOC}: case \texttt{a} uses the combined no-floc concentration,
while cases \texttt{b}/\texttt{c} use split \texttt{conc\_a}/\texttt{conc\_d}
fields.

\item \textbf{The effect is not active aggregation, bed exchange, or clipping.}
In \texttt{v006}, the aggregated component is zero throughout cases
\texttt{b}/\texttt{c}, the bed mass is zero, deposition/erosion exchange is
zero, split-minus-combined mismatch is zero, and cap counters are zero.  The
oscillatory mass loss therefore occurs while the floc branch is effectively a
disaggregated-only split transport problem.

\item \textbf{\texttt{conc\_d} is not numerically identical to \texttt{conc}.}
The physics expectation is correct: with \texttt{conc\_a=0} and no kinetic
transfer, \texttt{conc\_d} should behave like the no-floc combined
\texttt{conc}.  The implementation does not enforce that equivalence.  In
case \texttt{b}, \texttt{conc\_d} passes through the split-only path, including
floc-rate calculations at \texttt{mod\_plast.F:1021--1044}, split settling at
\texttt{mod\_plast.F:2729--2820}, split erosion/surface-flux updates at
\texttt{mod\_plast.F:2383--2440}, split OBC/source calls at
\texttt{mod\_plast.F:1543--1605}, and split clamp/exchange/recombine logic at
\texttt{mod\_plast.F:1657--1716}.  Case \texttt{a} uses the combined branches
for the same conceptual operators.  Therefore the \texttt{v006} difference can
arise from numerical/operator asymmetry even when aggregate physics is inactive.

\item \textbf{The day-3 timing points toward a transport or boundary trigger.}
Cases \texttt{a}, \texttt{b}, and \texttt{c} are indistinguishable through the
first two days and begin separating just before day 3.  After that, the no-floc
case increases smoothly, while the split cases alternate positive and negative
mass tendencies.  The most likely interpretation is that the split branch begins
interacting differently with a transport feature, open-boundary condition,
wet/dry edge, or strong-gradient region once the plastic plume reaches it.

\item \textbf{BBL roughness stays demoted.} The \texttt{v003} BBL-off run was
only \(9.13~\mathrm{kg}\) lower than the \texttt{v002} BBL-on run, while both
were more than \(410~\mathrm{kg}\) lower than \texttt{v005}.  BBL roughness is
not the leading explanation for the present case-\texttt{a} mass shift.

\item \textbf{Primary hypothetical cause.} The strongest current hypothesis is
that the \texttt{PLAST\_FLOC=T} split path is not algebraically equivalent to
the no-floc combined path when \texttt{SOURCE\_AG\_RATE=0},
\texttt{OBC\_AG\_RATE=0}, and aggregation remains zero.  The problem may be in
one of five nearby places: floc-rate setup before the normal update, split
deposition/erosion handling, the extra zero-field \texttt{adv\_scal}/
\texttt{vdif\_scal} calls for \texttt{conc\_a}, split OBC/source handling, or
the final clamp/exchange/recombine order.  The oscillation makes the OBC/source
or halo-exchange stages especially plausible, but this still needs an
operator-by-operator mass checkpoint.
\end{enumerate}

\section*{Suggested Next Move}
Keep \texttt{v006} as the clean comparison baseline and do one short,
high-instrumentation rerun around the divergence window.  The goal is to find
the first operator where \texttt{b} departs from \texttt{a}.

\begin{enumerate}
\item \textbf{Add stage mass checkpoints for case \texttt{a} and case
\texttt{b}.} For days \(2.75\)--\(3.25\), write domain-integrated MP1 mass after
each plastic operator: start of \texttt{Advance\_Plast}, after deposition,
after erosion/update-bottom, after horizontal advection, after vertical
diffusion, after OBC, after point-source BC, after nonnegative clamp, after MPI
exchange, and after recombination.  For \texttt{b}, log \texttt{conc\_a},
\texttt{conc\_d}, and \(\texttt{conc\_a}+\texttt{conc\_d}\) separately.

\item \textbf{Run a null-split equivalence test.} Add a temporary diagnostic
switch so \texttt{PLAST\_FLOC=T} allocates and outputs split fields, but skips
all aggregate work and transports only \texttt{conc\_d} through the exact
combined no-floc operator.  This means no floc-rate setup, no zero
\texttt{conc\_a} transport, no split OBC/source calls, no split cap, and no
split exchange/recombine order.  If this matches case \texttt{a}, the remaining
problem is inside the split operator sequence rather than setup or output.

\item \textbf{Bisect the split operator sequence.} Starting from the current
\texttt{b} setup, test these small changes one at a time: skip the zero
\texttt{conc\_a} \texttt{adv\_scal}/\texttt{vdif\_scal} calls when aggregate
rates are zero; apply OBC/source only to \texttt{conc\_d}; then force the
floc-branch final clamp/exchange/recombine order to mimic the no-floc branch as
closely as possible.

\item \textbf{Add an open-boundary/source budget.} Accumulate MP1 source input
and open-boundary export for the combined field in case \texttt{a} and for
\(\texttt{conc\_a}+\texttt{conc\_d}\) in case \texttt{b}.  If the day-3 split
coincides with a jump in export, the OBC path becomes the leading target.

\item \textbf{Check serial versus parallel behavior.} If the oscillation changes
with processor count, focus immediately on the split exchange calls around
\texttt{mod\_plast.F:1512--1516} and \texttt{mod\_plast.F:1679--1716}.
\end{enumerate}

\end{document}
